\subsection{Respecto a  algoritmos de Multiplicacion de Matrices}
El análisis experimental demuestra que una complejidad asintótica teórica superior no garantiza velocidad inmediata en la práctica debido al impacto de las constantes ocultas. En la multiplicación de matrices, la sobrecarga recursiva de Strassen satura el procesador tempranamente, permitiendo que las operaciones directas del método Naive sean más veloces para la dimensión máxima exigida de $N=2^8$. No obstante, la regla del límite asintótico establece que Strassen dominará inevitablemente a medida que el tamaño de entrada tienda al infinito.

\subsection{Respecto a algoritmos de Ordenamiento}

Por su parte, en los algoritmos de ordenamiento, el rendimiento depende críticamente de la gestión de memoria física y la mitigación de peores casos. QuickSort supera empíricamente a MergeSort porque opera in-place, evitando las costosas 2n asignaciones adicionales de memoria por nivel de recursión que MergeSort exige para mover datos. Además, reestructurar la lógica de los algoritmos previene su colapso computacional: centrar el pivote en QuickSort fuerza un árbol balanceado que evita su degradación a $O(n^2)$, mientras que la integración de un Min-Heap en PatienceSort garantiza la extracción del mínimo en un tiempo óptimo de O(1). Finalmente, la inclusión del algoritmo SORT de la librería estándar de C++ sirvió como el punto de referencia definitivo dentro de las mediciones experimentales, demostrando empíricamente la brecha de rendimiento que existe entre una implementación estandarizada y los algoritmos construidos manualmente.